\documentclass[12pt]{article}
\usepackage{amsmath}
\pagestyle{empty} % to turn off the page number

\begin{document}
Common Mathimatical Notation\\
Superscripts
$$2x^3$$
$$2x^34$$
$$2x^{34}$$
$$2x^{3x+4}$$
$$2x^{3x^4+5}$$
Subscripts
$$x_1$$
$$x_12$$
$$x_{12}$$
% $$x_1_2$$ error
$$x_{1_2}$$
$$x_{1_{2_3}}$$
$$ a_0, a_1, \ldots, a_{100}$$
Greek letters
$$\pi$$
$$\Pi$$
$$\alpha$$
$$ A=\pi r^2 $$
Trig functions 
$$ y=\sin x $$
$$ y=\cos x $$
$$ y=\csc \theta $$
$$ y=\sin^{-1} x $$
$$ y=\arcsin x $$
Log functions
$$y=\log x$$
$$y=\log_5 x$$
$$y=\ln x$$
Roots
$$\sqrt {2} $$
$$\sqrt [3]{2} $$
$$\sqrt{x^2+y^2}$$
$$\sqrt{1+\sqrt{x}}$$
Fractions
$$\frac{2}{3}$$
about $\frac{2}{3}$ of the glass is full.\\[6pt] % [6pt] to add some more vertical spacing
(you can make the fraction big like that : $\displaystyle \frac{2}{3}$ )\\[10pt]
Or it easy to type it like this : $\dfrac{2}{3}$ but you have to \textbackslash usepackage\{amsmath\}.
$$\frac{\sqrt{x+1}}{\sqrt{x+2}}$$
$$\frac{1}{1+\frac{1}{x}}$$
\end{document}